{\large\textbf{Mechanics of a Deformable Lattice}\hfill\textbf{(Total Marks : 20)}}

Here we study a deformable lattice hanging in gravity which acts as a deformable physical pendulum. It has only one degree of freedom, i.e.\ only one way to deform it and the configuration is fully described by an angle $\alpha$. Such structures have been studied by famous physicist James Maxwell in 19\textsuperscript{th} century, and some surprising behaviors have been discovered recently.

As shown in the figure 1, $N^2$ identical triangular plates (red triangle) are freely hinged by identical rods and form an $N \times N$ lattice ($N > 1$). The joints at the vertices are denoted by small circles. The sides of the equilateral triangles and the rods have the same length $l$. The dashed lines in the figure represent four tubes; each tube confines $N$ vertices (grey circles) on the edge and the $N$ vertices can slide in the tube, i.e.\ the tube is like a sliding rail.

The four tubes are connected in a diamond shape with two angles fixed at $60^\circ$ and another two angles at $120^\circ$ as shown in Figure 1. Each plate has a uniform density with mass $M$, and the other parts of the system are massless. The configuration of the lattice is uniquely determined by the angle $\alpha$, where $0^\circ \leq \alpha \leq 60^\circ$ (please see the examples of different angle $\alpha$ in Figure 1). The system is hung vertically like a “curtain” with the top tube fixed along the horizontal direction.

The coordinate system is shown in Figure 2. The zero level of the potential energy is defined at $y = 0$. A triangular plate is denoted by a pair of indices $(m,n)$, where $m, n = 0, 1, 2, \cdots, N - 1$ representing the order in the $x$ and $y$ directions respectively. A$(m,n)$, B$(m,n)$ and C$(m,n)$ denote the positions of the 3 vertices of the triangle $(m,n)$. The top-left vertex, A(0, 0) (the big black circle), is fixed.

The motion of the whole system is confined in the $x$-$y$ plane. The moment of inertia of a uniform equilateral triangular plate about its center of mass is $I = Ml^2/12$. The free fall acceleration is $g$. Please use $E_\mathrm{k}$ and $E_\mathrm{p}$ to denote kinetic energy and potential energy respectively.

\begin{center}
\includegraphics[width=1.00\linewidth]{figures/APhO_2016_Q1_fig1.png}

\textbf{Figure 1}
\end{center}

\begin{center}
\includegraphics[width=1.00\linewidth]{figures/APhO_2016_Q1_fig2.png}

\textbf{Figure 2}
\end{center}

\textbf{Section A:} When $N$=2 (as shown in figure 3):

\begin{center}
\includegraphics[width=1.00\linewidth]{figures/APhO_2016_Q1_fig3.png}

\textbf{Figure 3}
\end{center}

\noindent\begin{tabular}{|c|p{0.68\linewidth}|c|}
\hline
\textbf{A1} & What is the potential energy $E_\mathrm{p}$ of the system for a general angle $\alpha$ when $N = 2$? & \textbf{2 points} \\
\hline
\end{tabular}

\noindent\begin{tabular}{|c|p{0.68\linewidth}|c|}
\hline
\textbf{A2} & What is the equilibrium angle $\alpha_\mathrm{E}$ of the system under gravity when $N = 2$? & \textbf{1 point} \\
\hline
\end{tabular}

\noindent\begin{tabular}{|c|p{0.68\linewidth}|c|}
\hline
\textbf{A3} & The system follows a simple harmonic oscillation under a small perturbation from equilibrium. Calculate the kinetic energy of this system in terms of $\Delta\dot{\upalpha} \equiv \mathrm{d}(\Delta\upalpha)/\mathrm{d}t$. Calculate the oscillation frequency $f_\mathrm{E}$ when $N = 2$. & \textbf{5 points} \\
\hline
\end{tabular}

\textbf{Section B:} For arbitrary $N$ :

\noindent\begin{tabular}{|c|p{0.68\linewidth}|c|}
\hline
\textbf{B1} & What is the equilibrium angle $\alpha'_\mathrm{E}$ under gravity when $N$ is arbitrary? & \textbf{3 points} \\
\hline
\end{tabular}

\noindent\begin{tabular}{|c|p{0.68\linewidth}|c|}
\hline
\textbf{B2} & Consider the case when $N \rightarrow \infty$. Under a small perturbation of angle $\alpha$, the change of potential energy of the system is $\Delta E_\mathrm{p} \propto N^{\gamma_1}$, the kinetic energy of the system is $E_\mathrm{k} \propto N^{\gamma_2}$, and the oscillation frequency is $f'_\mathrm{E} \propto N^{\gamma_3}$. Find the values of $\gamma_1$, $\gamma_2$ and $\gamma_3$. & \textbf{3 points} \\
\hline
\end{tabular}

\textbf{Section C:} A force is exerted on one of the $3N^2$ triangle vertices so that the system maintains at $\upalpha_\mathrm{m} = 60^\circ$.

\begin{center}
\includegraphics[width=1.00\linewidth]{figures/APhO_2016_Q1_fig4.png}

\textbf{Figure 4}
\end{center}

\noindent\begin{tabular}{|c|p{0.68\linewidth}|c|}
\hline
\textbf{C1} & Which vertex should we choose to minimize the magnitude of this force? & \textbf{1 point} \\
\hline
\end{tabular}

\noindent\begin{tabular}{|c|p{0.68\linewidth}|c|}
\hline
\textbf{C2} & What are the direction and magnitude of this minimum force? Describe the direction in terms of the angle $\theta_F$ defined in Figure 4. & \textbf{5 points} \\
\hline
\end{tabular}
